\chapter[Band 9: Auswahlprinzip]{Band 9 auf einen Blick}
\noindent\textbf{Auswahlprinzip}\par\medskip
\noindent\emph{Lesart.} Die Surjektions-, Relations- und Familienform
des Auswahlprinzips werden miteinander verknüpft. Die Aussagen über
beliebige gleichzeitige Auswahlen verwenden ausdrücklich das Auswahlaxiom;
die Injektivität einer bereits vorhandenen Rechtsinversen benötigt es nicht.
\begin{BandUebersicht}
\textbf{Rechtsinverse}
Für \(F\colon A\to B\), \(G\colon B\to A\) bedeutet
\UebFormel{F\circ G=\Id_B,}
dass \(G\) eine Rechtsinverse von \(F\) ist.
&
\textit{Eine Rechtsinverse ist injektiv}
\UebFormel{F\circ G=\Id_B\vdash G\colon B\inj A.}
\UebQuellen{\FormulaRefAuto{F\circ G=\Id_B\vdash G\colon B\inj A}[thm]}\\
\addlinespace[.8ex]
\textbf{Auswahlaxiom in Surjektionsform}
Für jede Surjektion \(G\colon B\sur A\) wird gefordert:
\UebFormel{\exists F\colon A\to B\;G\circ F=\Id_A.}
\(\ACsur\) bezeichnet die entsprechende global quantifizierte Aussage.
\UebQuellen{\FormulaRefAuto{Auswahlaxiom}[ax];
\FormulaRefAuto{ACsur}[def]}
&
\textit{Auswahl aus einer totalen Relation}
Unter AC und \(\TotRel{R,A,B}\):
\UebFormel{\begin{aligned}
&\exists F\colon A\to B\,\forall x\in A\\
&\hspace{12mm}((x,F(x))\in R).
\end{aligned}}
\UebQuellen{\FormulaRefAuto{\exists F\colon A \to B\,\forall x\in A\,\bigl((x,F(x))\in R\bigr)}[thm]}\\
\addlinespace[.8ex]
\textbf{Disjunkte Familien und Auswahlmengen}
Für \(\DisjFam(M)\) soll \(L\) jedes Glied von \(M\) genau einmal
treffen. Die Mitgliedschaftsrelation aus Band 4 vermittelt zwischen
Familien- und Relationsform.
\UebQuellen{\FormulaRefAuto{Disjunkte Familie}[def]}
&
\textit{Familienform des Auswahlprinzips}
Unter AC und \(\DisjFam(M)\):
\UebFormel{\exists L\,\forall X\in M\,\exists!a\in X\,(a\in L).}
Der Band beweist außerdem die Rückrichtung zur Surjektionsform.
\UebQuellen{\FormulaRefAuto{\exists L\,\forall X\in M\,\exists! a\in X\,\bigl(a\in L\bigr)}[thm];
\FormulaRefAuto{\forall M\,\bigl(\DisjFam(M)\rightarrow \exists L\,\forall X\in M\,\exists! a\,(a\in X \land a\in L)\bigr) \vdash \exists F\colon A\to B\,\bigl(G\circ F=\Id_A\bigr)}[thm]}\\
\addlinespace[.8ex]
\textbf{Sektion aus einer Auswahlmenge}
Sei \(G\colon B\sur A\); \(L\) treffe jede Faser
\(X\in\Fib_G[A]\) genau einmal. Dann:
\UebFormel{\Sec_{G,L}(x)\coloneqq
\iota b\,(b\in\Fib_G(x)\land b\in L).}
\UebQuellen{\FormulaRefAuto{ChoiceSectionDef}[def]}
&
\textit{Die gewählte Sektion liefert eine Rechtsinverse}
\UebFormel{\begin{aligned}
&\forall X\in\Fib_G[A]\,\exists!a\,(a\in X\land a\in L)\\
&\hspace{6mm}\vdash\exists F\colon A\to B\,(G\circ F=\Id_A).
\end{aligned}}
\UebQuellen{\FormulaRefAuto{ChoiceSetRightInverse}[thm]}\\
\addlinespace[.8ex]
\textbf{Surjektive Größenvergleiche}
Die Surjektionsform von AC garantiert die Existenz einer Sektion;
deren Injektivität wurde am Anfang des Bandes unabhängig gezeigt.
&
\textit{Injektive Rechtsinverse}
Unter AC:
\UebFormel{\begin{aligned}
&F\colon A\sur B\vdash\exists H\colon B\inj A\\
&\hspace{18mm}\forall y\in B\,(F(H(y))=y).
\end{aligned}}
\UebQuellen{\FormulaRefAuto{F\colon A\sur B\vdash \exists H\colon B\inj A\,\forall y\in B\,\bigl(F(H(y))=y\bigr)}[thm]}\\
\addlinespace[.8ex]
\textbf{Retraktionen und ihre Fasern}
Seien \(G\subseteq S\), \(H\subseteq T\),
\(r\colon S\to G\), \(s\colon T\to H\) mit
\(r(g)=g\), \(s(h)=h\), und \(\varphi\colon G\bij H\).
\UebFormel{Q_g\coloneqq r^{-1}[\{g\}],\qquad
R_h\coloneqq s^{-1}[\{h\}].}
&
\textit{Bijektives Verkleben der Retraktionsfasern}
Unter AC und den links genannten Voraussetzungen:
\UebFormel{\begin{aligned}
&\forall g\in G\,\exists u\,(u\colon Q_g\bij R_{\varphi(g)})\\
&\vdash\exists f\,\bigl(f\colon S\bij T\\
&\quad\land\forall x\in S\,s(f(x))=\varphi(r(x))\\
&\quad\land\forall g\in G\,f(g)=\varphi(g)\bigr).
\end{aligned}}
\UebQuellen{\FormulaRefAuto{RetractionFiberBijection}[thm]}\\
\end{BandUebersicht}
